% Proof inputs: C_F0079 v5, including its fixed v4 interfaces; C_F0068 Part A for maximal-Delta rationality.
\section{Liberation dual system criterion}
\label{sec:liberation-dual-system-criterion}

We prove Theorem~\ref{thm:intrinsic-analytic} in this section.
We begin by recalling Definition~\ref{def:liberation-dual-system}.

\begin{definition*}
Let $X=(X_1,\ldots,X_d)\in\M^d$, and put $H:=L^2(\M,\tau)$ and $\Omega:=1$.
A \emph{liberation dual system} for $X$ is a tuple $(L_1,\ldots,L_d)$ of linear operators on $H$ such that
\[
 \M\Omega\subseteq\dom L_i\cap\dom L_i^*
 \qquad(1\leq i\leq d),
\]
and
\[
 [X_j,L_i]\xi=\delta_{ij}[X_i,P_\Omega]\xi
 \qquad(\xi\in\M\Omega,\ 1\leq i,j\leq d).
\]
Here $\M\Omega:=\{x\Omega:x\in\M\}$ and $P_\Omega$ is the orthogonal projection onto $\C\Omega$.
We call this system \emph{normalized} if
\[
 \sum_{i=1}^dL_i\xi=P_\Omega\xi
 \qquad(\xi\in\M\Omega).
\]
\end{definition*}

% Proof input: C_F0082, including both actual domains and sum normalization.
\begin{remark}\label{rem:liberation-gradients}
For $X\in\M^d$, put $A_i:=W^*(X_i)$ and $B_i:=W^*(X_j:j\ne i)$, with $B_1:=\C1$ if $d=1$.
If the liberation gradients $j(A_i:B_i)$ in the sense of \cite[Sections~5.3--5.4]{Voi99} exist in $L^2(\M,\tau)$ for every $i$, then $X$ admits a normalized liberation dual system in the sense of Definition~\ref{def:liberation-dual-system}.
\end{remark}

For a liberation dual system $(L_1,\ldots,L_d)$ and $\ell\in\M\otimes_{\C}\M^{\mathrm{op}}$, let $m_{L_i}(\ell):\M\Omega\to H$ be the linear operator defined by linear extension of
\[
 m_{L_i}(c\otimes d^{\mathrm{op}}):=cL_id|_{\M\Omega},\qquad c,d\in\M.
\]
The liberation commutator relations give, for $\ell\in\M\otimes_{\C}\M^{\mathrm{op}}$,
\begin{equation}\label{eq:liberation-tensor-actions}
 m_{L_i}(\Psi_j(\ell))=\delta_{ij}\Theta(\Psi_i(\ell))
 \quad\text{on }\M\Omega.
\end{equation}
Indeed, for $\ell=c\otimes d^{\mathrm{op}}$ with $c,d\in\M$, we have, on $\M\Omega$,
\[
 \begin{aligned}
 m_{L_i}(\Psi_j(\ell))
 &=m_{L_i}(cX_j\otimes d^{\mathrm{op}}-c\otimes(X_jd)^{\mathrm{op}})\\
 &=cX_jL_id-cL_iX_jd\\
 &=c[X_j,L_i]d\\
 &=\delta_{ij}c[X_i,P_\Omega]d\\
 &=\delta_{ij}\Theta(\Psi_i(\ell)).
 \end{aligned}
\]
The last equality uses equation~\eqref{eq:theta-bounded}.
Linearity gives equation~\eqref{eq:liberation-tensor-actions} for every $\ell$.

For $\ell\in\M\otimes_{\C}\M^{\mathrm{op}}$ and $s,t\in\M$, the definition of $m_{L_i}$ also gives
\begin{equation}\label{eq:liberation-phi-bimodule}
 m_{L_i}((s\otimes t^{\mathrm{op}})\ell)=sm_{L_i}(\ell)t
 \quad\text{on }\M\Omega.
\end{equation}
For $\ell=c\otimes d^{\mathrm{op}}$, multiplication in the opposite factor gives
\[
 (s\otimes t^{\mathrm{op}})(c\otimes d^{\mathrm{op}})
 =sc\otimes(dt)^{\mathrm{op}},
\]
and hence, on $\M\Omega$,
\[
 m_{L_i}((s\otimes t^{\mathrm{op}})(c\otimes d^{\mathrm{op}}))
 =scL_idt
 =sm_{L_i}(c\otimes d^{\mathrm{op}})t.
\]
Linearity proves equation~\eqref{eq:liberation-phi-bimodule} for every $\ell$.

\begin{proposition}\label{prop:liberation-normalization}
If $X\in\M^d$ admits a liberation dual system, then it admits a normalized liberation dual system.
\end{proposition}
\begin{proof}
Let $(L_1,\ldots,L_d)$ be a liberation dual system for $X$.
Define operators $C,\widetilde L_1,\ldots,\widetilde L_d$, each with domain $\M\Omega$, by
\[
 C\xi:=\sum_{k=1}^dL_k\xi-P_\Omega\xi,
 \qquad \widetilde L_i\xi:=L_i\xi-\frac1dC\xi
 \qquad(\xi\in\M\Omega,\ 1\leq i\leq d).
\]
Since $X_j\M\Omega\subseteq\M\Omega$, the liberation relations give, on $\M\Omega$,
\[
 [X_j,C]
 =\sum_{k=1}^d[X_j,L_k]-[X_j,P_\Omega]
 =\sum_{k=1}^d\delta_{kj}[X_k,P_\Omega]-[X_j,P_\Omega]
 =0.
\]
Consequently, for $\xi\in\M\Omega$ and $1\leq i,j\leq d$,
\[
 [X_j,\widetilde L_i]\xi
 =[X_j,L_i]\xi=\delta_{ij}[X_i,P_\Omega]\xi,
 \qquad
 \sum_{i=1}^d\widetilde L_i\xi
 =\sum_{i=1}^dL_i\xi-C\xi
 =P_\Omega\xi.
\]
Since $\mathcal M\Omega\subseteq\operatorname{dom}L_i^*$, we have for $\eta,\xi\in\M\Omega$,
\[
 \langle\eta,\widetilde L_i\xi\rangle
 =\left\langle
 L_i^*\eta-\frac1d\sum_{k=1}^dL_k^*\eta+\frac1dP_\Omega\eta,
 \xi\right\rangle.
\]
By the definition of the adjoint, \(\eta\in\operatorname{dom}\widetilde L_i^*\), and hence \(\mathcal M\Omega\subseteq\operatorname{dom}\widetilde L_i^*\).
Thus \(\widetilde L_1,\ldots,\widetilde L_d\) satisfy all the defining conditions.
\end{proof}

\begin{lemma}[{see \cite[Problem~III.5.18]{Kat95}}]\label{lem:finite-closure}
Let \(T\) be a densely defined closable operator on a Hilbert space \(H\).
If \(T\) has finite rank, then \(\overline T\in\F(H)\), with the same range and rank as \(T\).
\end{lemma}

\begin{lemma}\label{lem:tensor-difference-injectivity}
Let $V,W$ be complex vector spaces, and let $A:V\to V$ and $B:W\to W$ be linear maps.
If $A$ or $B$ has no nonzero finite-dimensional invariant subspace, then
\[
 A\otimes I_W-I_V\otimes B:V\otimes_{\C}W\longrightarrow V\otimes_{\C}W
\]
is injective.
\end{lemma}
\begin{proof}
Suppose that $B$ has no nonzero finite-dimensional invariant subspace.
If a nonzero kernel vector existed, write it as $\sum_{r=1}^n v_r\otimes w_r$ with $n$ minimal.
By minimality of $n$, both families are linearly independent.
Set $E:=\operatorname{span}\{w_1,\ldots,w_n\}$, and let $\pi:W\to W/E$ be the quotient map.
Since the tensor lies in the kernel, we have
\[
 \sum_{r=1}^n Av_r\otimes w_r=\sum_{r=1}^n v_r\otimes Bw_r.
\]
Applying $I_V\otimes \pi$ and using $\pi(w_r)=0$ gives
\[
 \sum_{r=1}^n v_r\otimes \pi(Bw_r)=0.
\]
Since $v_1,\ldots,v_n$ are linearly independent, choose linear functionals $\varphi_j:V\to\C$ satisfying $\varphi_j(v_r)=\delta_{jr}$.
Applying $\varphi_j\otimes I_{W/E}$ yields $\pi(Bw_j)=0$, and hence $Bw_j\in E$, for every $1\leq j\leq n$.
Thus $B(E)\subseteq E$, contradicting the assumption on $B$, since $E$ is nonzero and finite-dimensional.
The case where $A$ has the stated property follows by interchanging the tensor factors.
\end{proof}

\begin{lemma}\label{lem:tensor-quotient-kernels}
Let $E$ be a subspace of a complex vector space $V$, and let $\pi:V\to V/E$ be the quotient map.
For every complex vector space $W$,
\[
 \ker(\pi\otimes I_W)=E\otimes_{\C}W,
 \qquad
 \ker(I_W\otimes\pi)=W\otimes_{\C}E.
\]
\end{lemma}
\begin{proof}
Choose a complement $V_1$ such that $V=E\oplus V_1$.
Then
\[
 V\otimes_{\C}W
 =(E\otimes_{\C}W)\oplus(V_1\otimes_{\C}W),
\]
and the map $\pi\otimes I_W$ vanishes on the first summand and maps the second summand isomorphically onto $(V/E)\otimes_{\C}W$, since $\pi|_{V_1}:V_1\to V/E$ is an isomorphism.
This proves the first kernel equality.
Interchanging the tensor factors proves injectivity of the other natural map and the second kernel equality.
\end{proof}

\begin{lemma}\label{lem:tensor-l2-recovery}
Let $x\in\M$ and suppose that $x-\lambda$ is regular for every $\lambda\in\C$.
The map $\Psi_x$ of Definition~\ref{def:psi} is injective and satisfies
\[
 \Psi_x^{-1}(H\otimes_{\C}H)=H\otimes_{\C}H.
\]
\end{lemma}
\begin{proof}
Put $U:=L^0(\M)$ and define the right multiplication maps
\[
 \begin{aligned}
 R_x:U&\longrightarrow U,& R_x(u)&:=ux,\\
 R_{x^{\mathrm{op}}}:U^{\mathrm{op}}&\longrightarrow U^{\mathrm{op}},&
 R_{x^{\mathrm{op}}}(v^{\mathrm{op}})&:=(xv)^{\mathrm{op}}.
 \end{aligned}
\]
Thus $\Psi_x=R_x\otimes I-I\otimes R_{x^{\mathrm{op}}}$.
Neither right multiplication map has a nonzero finite-dimensional invariant subspace: such a subspace would contain an eigenvector, contradicting regularity of $x-\lambda$.
Lemma~\ref{lem:tensor-difference-injectivity} therefore implies that $\Psi_x$ is injective.

For the preimage equality, observe that $x\in\M$ implies that both maps preserve $H$.
Their restrictions to $H$ also have no nonzero finite-dimensional invariant subspace.
In particular, $\Psi_x$ preserves $H\otimes_{\C}H$.

Let $\pi_1:U\to U/H$ and $\pi_2:U^{\mathrm{op}}\to U^{\mathrm{op}}/H$ be the quotient maps.
The induced linear maps on the quotient spaces are
\[
 \begin{aligned}
 \widetilde R_x:U/H&\longrightarrow U/H,&
 \widetilde R_x(u+H)&:=ux+H,\\
 \widetilde R_{x^{\mathrm{op}}}:U^{\mathrm{op}}/H&\longrightarrow U^{\mathrm{op}}/H,&
 \widetilde R_{x^{\mathrm{op}}}(v^{\mathrm{op}}+H)&:=(xv)^{\mathrm{op}}+H.
 \end{aligned}
\]
They satisfy
\[
 \pi_1R_x=\widetilde R_x\pi_1,
 \qquad
 \pi_2R_{x^{\mathrm{op}}}=\widetilde R_{x^{\mathrm{op}}}\pi_2.
\]
Now let $k\in U\otimes_{\C}U^{\mathrm{op}}$ satisfy $\Psi_x(k)\in H\otimes_{\C}H$.
Using $\pi_1R_x=\widetilde R_x\pi_1$, we obtain
\[
 \begin{aligned}
 &\bigl(\widetilde R_x\otimes I-I\otimes R_{x^{\mathrm{op}}}\bigr)
    (\pi_1\otimes I)k\\
 &\qquad=(\pi_1\otimes I)\Psi_x(k)=0.
 \end{aligned}
\]
The last equality follows from $\Psi_x(k)\in H\otimes_{\C}H$ and $\pi_1|_H=0$.
By Lemma~\ref{lem:tensor-difference-injectivity}, the map $\widetilde R_x\otimes I-I\otimes R_{x^{\mathrm{op}}}$ is injective on $(U/H)\otimes_{\C}U^{\mathrm{op}}$, since $R_{x^{\mathrm{op}}}$ has no nonzero finite-dimensional invariant subspace.
Consequently, $(\pi_1\otimes I)k=0$.
Applying the first kernel equality of Lemma~\ref{lem:tensor-quotient-kernels} with $V=U$, $E=H$, $W=U^{\mathrm{op}}$ and $\pi=\pi_1$ gives
\[
 k\in\ker(\pi_1\otimes I)=H\otimes_{\C}U^{\mathrm{op}}.
\]

Now regard $k$ as an element of $H\otimes_{\C}U^{\mathrm{op}}$.
Using $\pi_2R_{x^{\mathrm{op}}}=\widetilde R_{x^{\mathrm{op}}}\pi_2$, we similarly obtain
\[
 k\in\ker(I\otimes \pi_2)=H\otimes_{\C}H.
\]
\end{proof}

\begin{proposition}\label{prop:liberation-delta}
Let $X=(X_1,\ldots,X_d)\in\M^d$, and suppose that $X_i-\lambda$ is regular for every $1\leq i\leq d$ and every $\lambda\in\C$.
If $X$ admits a liberation dual system, then $\Delta(X)=d$.
\end{proposition}
\begin{proof}
Let $T_1,\ldots,T_d\in\F(H)$ satisfy $\sum_i[\rho(X_i),T_i]=0$, and put $k_i:=\Theta^{-1}(T_i)\in H\otimes_{\C}H$.
Since $X_i\in\M$, each $\Psi_i(k_i)$ belongs to $H\otimes_{\C}H$.
By equation~\eqref{eq:tensor-commutator-general},
\[
 \Theta\!\left(\sum_i\Psi_i(k_i)\right)
 =\sum_i\Theta(\Psi_i(k_i))
 =\sum_i[\rho(X_i),T_i]=0.
\]
Since $\Theta$ is injective, $\sum_i\Psi_i(k_i)=0$.
By Lemma~\ref{lem:tensor-clearing}, choose regular $s,t\in\M$ such that $\ell_i:=(s\otimes t^{\mathrm{op}})k_i\in\M\otimes_{\C}\M^{\mathrm{op}}$ for every $i$.
Equation~\eqref{eq:psi-left-linearity} yields $\sum_i\Psi_i(\ell_i)=0$.
Applying $m_{L_i}$ and equation~\eqref{eq:liberation-tensor-actions} gives, on $\M\Omega$,
\[
 0=m_{L_i}\!\left(\sum_j\Psi_j(\ell_j)\right)
 =\Theta(\Psi_i(\ell_i)).
\]
Since $\Theta(\Psi_i(\ell_i))$ is bounded and $\M\Omega$ is dense in $H$, it vanishes on $H$.
Hence $\Psi_i(\ell_i)=0$.
Since $X_i-\lambda$ is regular for every $\lambda\in\C$, Lemma~\ref{lem:tensor-l2-recovery} gives $\ell_i=0$.
Cancellation of $s\otimes t^{\mathrm{op}}$ gives $k_i=0$, and therefore $T_i=0$.
Remark~\ref{rem:cycle-characterization} now gives $\Delta(X)=d$.
\end{proof}

\begin{lemma}\label{lem:tensor-images-l2}
Let $X=(X_1,\ldots,X_d)\in\M^d$, and suppose that $X_i-\lambda$ is regular for every $1\leq i\leq d$ and every $\lambda\in\C$.
Let $(L_1,\ldots,L_d)$ be a liberation dual system for $X$.
If $k_1,\ldots,k_d\in L^0(\M)\otimes_{\C}L^0(\M)^{\mathrm{op}}$ satisfy
\[
 \ell_0:=\sum_{i=1}^d\Psi_i(k_i)\in\M\otimes_{\C}\M^{\mathrm{op}},
\]
then, for every $i$, $k_i\in H\otimes_{\C}H$ and
\begin{equation}\label{eq:liberation-tensor-closure}
 \overline{m_{L_i}(\ell_0)}=\Theta(\Psi_i(k_i)).
\end{equation}
\end{lemma}
\begin{proof}
Fix $i$, write $\ell_0=\sum_r c_r\otimes d_r^{\mathrm{op}}$ with $c_r,d_r\in\M$, and put $T_i:=m_{L_i}(\ell_0)$ on $\M\Omega$.
By Definition~\ref{def:liberation-dual-system}, $\M\Omega\subseteq\dom L_i^*$.
Since $c_r^*\M\Omega\subseteq\M\Omega$, the definition of the adjoint gives
\[
 \M\Omega\subseteq\dom T_i^*,\qquad
 T_i^*\eta=\sum_r d_r^*L_i^*(c_r^*\eta)
 \quad(\eta\in\M\Omega).
\]
Thus $T_i$ is closable.
By Lemma~\ref{lem:tensor-clearing}, choose regular $s,t\in\M$ such that
\[
 \ell_j:=(s\otimes t^{\mathrm{op}})k_j
 \in\M\otimes_{\C}\M^{\mathrm{op}}
 \qquad(1\leq j\leq d).
\]
Equation~\eqref{eq:psi-left-linearity} gives
\[
 \sum_j\Psi_j(\ell_j)=(s\otimes t^{\mathrm{op}})\ell_0.
\]
Using equation~\eqref{eq:liberation-phi-bimodule} in the second equality below and equation~\eqref{eq:liberation-tensor-actions} in the last equality, we obtain, on $\M\Omega$,
\begin{equation}\label{eq:liberation-cleared-operator}
 \begin{aligned}
 sT_it
 &=sm_{L_i}(\ell_0)t\\
 &=m_{L_i}((s\otimes t^{\mathrm{op}})\ell_0)\\
 &=m_{L_i}\!\left(\sum_j\Psi_j(\ell_j)\right)\\
 &=\Theta(\Psi_i(\ell_i)).
 \end{aligned}
\end{equation}
The operator $\Theta(\Psi_i(\ell_i))$ has finite rank and $s$ is injective, so the restriction of $T_i$ to $t\M\Omega$ has finite rank.
Since $t$ has dense range, $t\M\Omega$ is dense in $H$.
This restriction is closable, and Lemma~\ref{lem:finite-closure} makes its closure a bounded finite-rank operator on $H$.
The operator $\overline{T_i}$ extends that closure and therefore equals it.
Set $h_i:=\Theta^{-1}(\overline{T_i})\in H\otimes_{\C}H$.
Extending equation~\eqref{eq:liberation-cleared-operator} by boundedness and density, we have
\[
 s\overline{T_i}t=\Theta(\Psi_i(\ell_i)).
\]
Applying $\Theta^{-1}$ to both sides, we obtain
\[
 (s\otimes t^{\mathrm{op}})h_i
 =\Psi_i(\ell_i)
 =(s\otimes t^{\mathrm{op}})\Psi_i(k_i).
\]
Cancellation of the regular outside factors gives $\Psi_i(k_i)=h_i\in H\otimes_{\C}H$.
Since every $X_i-\lambda$ is regular, Lemma~\ref{lem:tensor-l2-recovery} gives $k_i\in H\otimes_{\C}H$.
\end{proof}

The proof of Theorem~\ref{thm:intrinsic-analytic} uses the following closability result for fraction commutators.

\begin{lemma}\label{lem:fraction-domains}
Let \(D_1,\ldots,D_d\) be densely defined linear operators on \(H\) such that
\[\M\Omega\subseteq\dom D_i\cap\dom D_i^*\]
for \(1\leq i\leq d\).
For every \(u\in L^0(\M)\), every fraction pair \(Q\) of \(u\), and every \(i\), the operator \(\Gamma[Q,D_i]\) is closable and densely defined.
\end{lemma}
\begin{proof}
Fix \(i\) and write \(Q=(s,a,b,t)\).
Since \(b,t\in\M\) preserve \(\M\Omega\), the domain of \(\Gamma[Q,D_i]\) contains \(\M\Omega\).
For \(\eta\in\M\Omega\), we have \(s^*\eta,a^*\eta\in\M\Omega\subseteq\dom D_i^*\).
Thus, for \(\xi\in\dom\Gamma[Q,D_i]\),
\begin{equation}\label{eq:defect-adjoint}
 \langle\eta,\Gamma[Q,D_i]\xi\rangle
 =\left\langle
 b^*D_i^*(s^*\eta)-t^*D_i^*(a^*\eta),\xi
 \right\rangle.
\end{equation}
Thus \(\M\Omega\subseteq\dom(\Gamma[Q,D_i])^*\), so \(\Gamma[Q,D_i]\) is closable.
\end{proof}

\begin{proof}[Proof of Theorem~\ref{thm:intrinsic-analytic}]
By Proposition~\ref{prop:liberation-delta}, $\Delta(X)=d$.
We prove $(1)\implies(4)\implies(3)\implies(2)\implies(1)$.

$(1)\implies(4)$.
Fix any fraction pair $Q=(s,a,b,t)$ for $u$.
By Theorem~\ref{thm:delta}\textup{(\ref{item:delta-all-fraction-tensors})}, there are tensors $\ell_i\in L^0(\M)\otimes_{\C}L^0(\M)^{\mathrm{op}}$ satisfying equation~\eqref{eq:all-fraction-tensor-criterion}.
Set
\[
 \ell_0:=a\otimes t^{\mathrm{op}}-s\otimes b^{\mathrm{op}}
   =\sum_i\Psi_i(\ell_i).
\]
Lemma~\ref{lem:tensor-images-l2} gives $\ell_i\in H\otimes_{\C}H$ and $\overline{m_{L_i}(\ell_0)}=\Theta(\Psi_i(\ell_i))$.
For $\xi\in\M\Omega$, by the definition of $m_{L_i}$ and $\ell_0$,
\[
 m_{L_i}(\ell_0)\xi
 = aL_i(t\xi)-sL_i(b\xi)
 = -\Gamma[Q,L_i]\xi.
\]
By Definition~\ref{def:liberation-dual-system}, $\M\Omega\subseteq\dom L_i\cap\dom L_i^*$.
Lemma~\ref{lem:fraction-domains} therefore implies that $\Gamma[Q,L_i]$ is closable.
Thus we have
\[
\overline{\Gamma[Q,L_i]}=-\overline{m_{L_i}(\ell_0)}=-\Theta(\Psi_i(\ell_i)).
\]
With $T_i:=-\Theta(\ell_i)\in\F(H)$, equation~\eqref{eq:tensor-commutator-general} now gives
\[
 \overline{\Gamma[Q,L_i]}
 =-\Theta(\Psi_i(\ell_i))
 =[\rho(X_i),T_i].
\]
Thus assertion~\textup{(4)} holds.

$(4)\implies(3)$.
Fix any fraction pair $Q=(s,a,b,t)$ for $u$ and choose $T_i\in\F(H)$ such that
$[\rho(X_i),T_i]=\overline{\Gamma[Q,L_i]}$ for every $i$.
By equation~\eqref{eq:liberation-normalization}, on $\M\Omega$ we have
\[
 \begin{aligned}
  \sum_i[\rho(X_i),T_i]
  &=\sum_i\overline{\Gamma[Q,L_i]}\\
  &=s\Bigl(\sum_iL_i\Bigr)b-a\Bigl(\sum_iL_i\Bigr)t\\
  &=sP_\Omega b-aP_\Omega t
   =\Gamma[Q,P_\Omega].
 \end{aligned}
\]
The bounded operators $\sum_i[\rho(X_i),T_i]$ and $\Gamma[Q,P_\Omega]$ agree on the dense subspace $\M\Omega$ and hence on all of $H$.
This proves assertion~\textup{(3)}.

$(3)\implies(2)$.
Let $Q$ be the fraction pair for $u$ given by equation~\eqref{eq:canonical}.
By assertion~\textup{(3)}, there exist $T_1,\ldots,T_d\in\F(H)$ satisfying equation~\eqref{eq:intrinsic-analytic-criterion} for this $Q$.

$(2)\implies(1)$.
This is the implication $\textup{(\ref{item:delta-finite-rank})}\implies\textup{(1)}$ of Theorem~\ref{thm:delta}.

For each fixed fraction pair, uniqueness of the tuple $(T_1,\ldots,T_d)$ follows from Theorem~\ref{thm:delta}.
\end{proof}
